\documentclass[11pt,a4paper]{article}

\usepackage[utf8]{inputenc}
\usepackage[T1]{fontenc}
\usepackage{amsmath}
\usepackage{amssymb}
\usepackage{booktabs}
\usepackage{array}
\usepackage[margin=2.5cm]{geometry}
\usepackage{fancyhdr}
\usepackage{hyperref}

% Config and file paths are long unbreakable typewriter strings. \nolinkurl
% lets them break at '.', '_' and '/' with no hyphen and no hyperlink; \texttt
% cannot break at all and runs past the margin. Pandoc renders both as code.
\urlstyle{tt}
\Urlmuskip=0mu plus 1mu
\def\UrlBreaks{\do\.\do\_\do\/\do\-\do\:\do\,}

\numberwithin{equation}{section}

\pagestyle{fancy}
\fancyhf{}
\lhead{bess-stack-de}
\rhead{V1.0}
\lfoot{Can Cokacar, M.Sc.}
\rfoot{\thepage}
\renewcommand{\headrulewidth}{0.4pt}
\renewcommand{\footrulewidth}{0.4pt}

\title{Mathematical Formulation of the Dispatch Optimization Algorithm}
\author{bess-stack-de \\ Degradation-aware rolling-horizon dispatch \\
for a grid-connected BESS in the German market}
\date{V1.0}

\begin{document}

\maketitle
\thispagestyle{fancy}

\section{Scope and Conventions}

This document states the optimization model implemented in
\nolinkurl{src/bess_stack/model/dispatch.py}. The project models

\begin{quote}
a degradation-aware, rolling-horizon dispatch optimizer for a single
grid-connected BESS in the German market, co-optimizing day-ahead arbitrage,
FCR capacity, and aFRR capacity + energy, that outputs an annual revenue stack
and feeds a project-level IRR calculation.
\end{quote}

\textbf{Section 2 describes what is implemented. Section 3 describes what is
not.} Every equation in Section 2 corresponds to a statement in
\nolinkurl{solve_window}; no equation there describes intended behaviour. Section 3
formulates the reserve products from the parameters already present in
\nolinkurl{scenarios/reference.yaml}, none of which the code reads today. The
distinction is maintained deliberately: a formulation that does not say which
half is running is not a specification but an advertisement.

\subsection{Units and sign conventions}

Power is in MW, energy in MWh, prices and costs in EUR/MWh, and capacity prices
in EUR/MW/h. One time step spans $\Delta t$ hours. Charging and discharging
power are separate non-negative variables rather than one signed variable, which
is what allows the complementarity condition of constraint group B1 to be written
as a linear constraint. State of charge is carried as an energy $E_t$ in MWh rather
than a percentage; the scenario file states its bounds as fractions and the
model multiplies them by the nominal capacity on load.

\section{Day-Ahead Arbitrage Model (Implemented)}

\subsection{Index sets}

\begin{tabular}{@{}l >{\raggedright\arraybackslash}p{12cm}@{}}
\toprule
Symbol & Description \\
\midrule
$t \in \mathcal{T}$ & Time steps within one optimization window,
                      $\mathcal{T} = \{0, 1, \dots, N-1\}$ \\
$\mathcal{B}_j$     & Steps belonging to day-ahead product block $j$, so that
                      $\bigcup_j \mathcal{B}_j = \mathcal{T}$ \\
\bottomrule
\end{tabular}

\subsection{Parameters}

Each parameter names the field it is read from, so that the document and the
configuration cannot drift apart unnoticed. Source fields are dotted paths into
\nolinkurl{scenarios/reference.yaml}; those beginning with a stream name
(\nolinkurl{day_ahead}, \texttt{fcr}, \texttt{afrr}) are relative to
\nolinkurl{revenue_streams}.

\small
\begin{tabular}{@{}l >{\raggedright\arraybackslash}p{4.4cm} l >{\raggedright\arraybackslash}p{5.6cm}@{}}
\toprule
Symbol & Description & Unit & Source field \\
\midrule
$\Delta t$   & Length of one time step & h & \nolinkurl{market.resolution_minutes} \\
$\pi_t$      & Day-ahead price in step $t$ & EUR/MWh & price series argument \\
$\kappa$     & Spread retention, $1 - \text{haircut}$ & -- & \nolinkurl{day_ahead.spread_haircut} \\
$f$          & Trading fee per MWh traded & EUR/MWh & \nolinkurl{day_ahead.fees_eur_per_mwh} \\
$g^{c}$      & Network fee on energy charged & EUR/MWh & \nolinkurl{grid.charging_fees_eur_per_mwh} \\
$g^{d}$      & Network fee on energy discharged & EUR/MWh & \nolinkurl{grid.discharging_fees_eur_per_mwh} \\
$c^{\deg}$   & Marginal degradation cost & EUR/MWh & \nolinkurl{degradation.marginal_cost_eur_per_mwh} \\
$\bar{P}$    & Binding power cap & MW & $\min$(\nolinkurl{battery.power_mw}, \nolinkurl{grid.connection_limit_mw}) \\
$E^{\text{nom}}$ & Nominal usable energy & MWh & \nolinkurl{battery.energy_mwh} \\
$\eta^{c}$   & One-way charging efficiency & -- & \nolinkurl{battery.efficiency.charge} \\
$\eta^{d}$   & One-way discharging efficiency & -- & \nolinkurl{battery.efficiency.discharge} \\
$\sigma$     & Self-discharge, fraction of $E^{\text{nom}}$ per day & 1/d & \nolinkurl{battery.efficiency.self_discharge_per_day} \\
$P^{\text{aux}}$ & Parasitic auxiliary load & MW & \nolinkurl{battery.aux_load_mw} \\
$L$          & Idle energy loss per step, see (2.5) & MWh & derived \\
$E^{\min}$   & Lower energy bound & MWh & \nolinkurl{battery.soc.min} $\times\, E^{\text{nom}}$ \\
$E^{\max}$   & Upper energy bound & MWh & \nolinkurl{battery.soc.max} $\times\, E^{\text{nom}}$ \\
$E_{\text{init}}$ & Energy at window start & MWh & carried between windows \\
$E^{\text{term}}$ & Required energy at window end & MWh & \nolinkurl{dispatch.terminal_soc_fraction} $\times\, E^{\text{nom}}$ \\
$m$          & Time steps per day-ahead block & -- & \nolinkurl{day_ahead.product_resolution_minutes} \\
$H$          & Optimization horizon & h & \nolinkurl{dispatch.horizon_hours} \\
$C$          & Committed step of the rolling horizon & h & \nolinkurl{dispatch.commit_step_hours} \\
\bottomrule
\end{tabular}
\normalsize

\subsection{Decision variables}

\small
\begin{tabular}{@{}l l l >{\raggedright\arraybackslash}p{7.5cm}@{}}
\toprule
Variable & Domain & Unit & Description \\
\midrule
$P^{c}_{t}$ & $[0, \bar{P}]$ & MW & Charging power in step $t$ \\
$P^{d}_{t}$ & $[0, \bar{P}]$ & MW & Discharging power in step $t$ \\
$u_{t}$     & $\{0,1\}$      & -- & Mode indicator, $1$ = charging permitted, $0$ = discharging permitted \\
$E_{t}$     & $[E^{\min}, E^{\max}]$ & MWh & Stored energy at the end of step $t$ \\
\bottomrule
\end{tabular}
\normalsize

\subsection{Objective function}

The model maximizes net trading revenue over the window:

\begin{equation}
\max \; \sum_{t \in \mathcal{T}}
  \Big[ \big( \pi_t \kappa - f - g^{d} - c^{\deg} \big) P^{d}_{t} \Delta t
      - \big( \pi_t \kappa + f + g^{c} \big) P^{c}_{t} \Delta t \Big]
\end{equation}

Equation (2.1) is the sum over all steps of revenue from selling less the cost
of buying. Its terms are:

\begin{itemize}
  \item $\pi_t \kappa$ --- the price actually captured. The haircut $\kappa$
        shaves the traded spread to represent the gap between the price used in
        the model and the price a real bid achieves.
  \item $f$, $g^{c}$, $g^{d}$ --- per-MWh fees. Note the asymmetry of sign: fees
        subtract from the sale price and \emph{add} to the purchase price, so
        they narrow the spread from both ends rather than acting as a levy on
        net position.
  \item $c^{\deg}$ --- the marginal cost of degradation, charged per MWh
        discharged. This is not a cash cost and nobody invoices it. It is the
        shadow price of consuming battery life, and it is what makes the model
        degradation-aware rather than merely degradation-reporting: a cycle
        whose spread does not cover $c^{\deg}$ is declined. Its admissible
        range is bounded above by the cost of making good the capacity one
        cycle consumes,
        \[
          c^{\deg} \le \phi \cdot \kappa^{\text{aug}} \cdot 1000 ,
        \]
        where $\phi$ is the cyclic fade per full cycle and
        $\kappa^{\text{aug}}$ the augmentation cost per kWh, both read from the
        \nolinkurl{degradation} block. Nominal energy
        cancels out of the ratio, because fade per cycle is a fraction of
        beginning-of-life capacity and one equivalent full cycle discharges
        exactly that. \nolinkurl{config.py} enforces the bound. There is no
        corresponding lower bound: establishing one needs a discount factor
        keyed to the augmentation date, which depends on the cycle count, which
        is an outcome of dispatch rather than an input to it. Satisfying the
        bound therefore shows the value is not indefensibly high, not that it is
        right.
  \item $\Delta t$ --- converts power in MW to energy in MWh over the step.
\end{itemize}

Because degradation is charged on discharge only, throughput is measured as
discharged energy throughout this document.

\subsection{Operational constraints}

\subsubsection*{B1 --- Charge/discharge complementarity and power limits}

A battery may not charge and discharge in the same step. The restriction is
enforced through the binary mode indicator, which simultaneously imposes the
power limit:

\begin{align}
P^{c}_{t} &\le \bar{P} \, u_{t} & \forall t \in \mathcal{T} \\
P^{d}_{t} &\le \bar{P} \, (1 - u_{t}) & \forall t \in \mathcal{T}
\end{align}

where $\bar{P}$ is the smaller of the battery power rating and the grid
connection limit, and $u_t$ selects the mode.

This binary is not a modelling nicety and must not be relaxed to obtain a linear
program. Under the relaxation, at sufficiently negative prices the optimizer
charges and discharges simultaneously: the two positions cancel in the objective
at zero net cost, while round-trip losses destroy stored energy. Destroying
energy for free is valuable when prices are negative, because it lets a full
battery keep buying. The behaviour is an artefact of the relaxation rather than
a dispatch strategy, and the binary is what forbids it.

\subsubsection*{B2 --- State-of-charge dynamics}

Stored energy evolves under charging, discharging and idle losses:

\begin{equation}
E_{t} = E_{t-1}
      + \eta^{c} P^{c}_{t} \Delta t
      - \frac{P^{d}_{t} \Delta t}{\eta^{d}}
      - L
      \qquad \forall t \in \mathcal{T}
\end{equation}

with $E_{-1} = E_{\text{init}}$, and the idle loss per step

\begin{equation}
L = \left( \frac{\sigma E^{\text{nom}}}{24} + P^{\text{aux}} \right) \Delta t .
\end{equation}

where $\eta^{c}$ \emph{multiplies} the charging term and $\eta^{d}$
\emph{divides} the discharging term. The asymmetry is deliberate and is the
correct treatment of one-way efficiencies: of the energy drawn from the grid
only a fraction $\eta^{c}$ is stored, whereas delivering $P^{d}_{t}$ to the grid
requires drawing $P^{d}_{t} / \eta^{d}$ from storage. The product
$\eta^{c} \eta^{d}$ is the round-trip efficiency. The loss term $L$ is
independent of dispatch and applies in every step, including idle ones.

\subsubsection*{B3 --- Energy bounds}

\begin{equation}
E^{\min} \le E_{t} \le E^{\max} \qquad \forall t \in \mathcal{T}
\end{equation}

where the bounds are the scenario's state-of-charge fractions scaled by nominal
capacity. In the implementation these are variable bounds rather than explicit
rows.

\subsubsection*{B4 --- Product block coherence}

Day-ahead energy settles in hourly blocks, so dispatch may not vary within a
block even when the model runs at a finer resolution:

\begin{align}
P^{c}_{t} &= P^{c}_{t-1} & \forall t : t \bmod m \ne 0 \\
P^{d}_{t} &= P^{d}_{t-1} & \forall t : t \bmod m \ne 0
\end{align}

where $m$ is the number of model steps in one product block. With a 15-minute
resolution and a 60-minute day-ahead product, $m = 4$. This is also the
machinery the four-hour reserve blocks of Section 3 will reuse, which is why it
is expressed generally rather than hard-coded to the hourly case.

\subsubsection*{B5 --- Terminal state of charge}

\begin{equation}
E_{N-1} \ge E^{\text{term}}
\end{equation}

where $E^{\text{term}}$ is the required closing energy. The constraint exists
because stored energy has no value beyond the horizon: without it the optimizer
empties the battery into the final step of every window, since selling is free
revenue and the consequence falls outside the model. It is applied at the far
end of the window, not at the end of the committed region, so that it shapes end
effects rather than the decisions actually taken.

\subsection{Rolling-horizon procedure}

The window model above is solved repeatedly across the price series. With
$h = H / \Delta t$ steps of horizon and $c = C / \Delta t$ steps of commitment:

\begin{enumerate}
  \item Set $E_{\text{init}}$ to the opening state of charge.
  \item Solve (2.1)--(2.9) over the next $h$ steps.
  \item Retain the first $c$ steps of the solution as committed dispatch.
  \item Set $E_{\text{init}}$ to $E_{c-1}$ from the solved window.
  \item Advance by $c$ steps and repeat until the series is exhausted.
\end{enumerate}

The horizon exceeds the commitment ($h > c$) so that the uncommitted tail
absorbs the distortion of B5. A final window shorter than one commitment step
receives no terminal target, since there is no tail left to absorb it.

Setting \nolinkurl{dispatch.mode} to \nolinkurl{perfect_foresight} collapses the
procedure to a single window spanning the entire series. This reports an upper
bound on achievable revenue rather than an expectation, and the same is true of
the rolling mode while \nolinkurl{dispatch.forecast.kind} remains \texttt{perfect}:
each window sees its own prices exactly.

\subsection{Post-optimization network charges}

Network charges that fall outside the exemption of \S118(6) EnWG are applied in
\nolinkurl{src/bess_stack/finance/grid_fees.py} to an already-solved dispatch,
rather than by re-solving each project year under that year's tariff. For a year
$y$ the charge is

\begin{equation}
\Phi_y = \big( \gamma_y - \gamma^{\text{priced}} \big) \sum_{t} P^{c}_{t} \Delta t
       + \mathbb{1}[\,y \text{ not exempt}\,] \cdot \rho \cdot \bar{P}^{\text{conn}} \cdot 1000
\end{equation}

where $\gamma_y$ is the per-MWh charging fee applicable in year $y$,
$\gamma^{\text{priced}}$ is the fee the optimizer already carried in (2.1) via
$g^{c}$, $\rho$ is the successor capacity charge in EUR/kW/year and
$\bar{P}^{\text{conn}}$ is the connection capacity in MW. Subtracting the fee
already priced is what prevents a charge from being billed twice.

This is an approximation and its error is measured rather than assumed. Because
the retained dispatch was optimized without the fee, it remains feasible but no
longer optimal once the fee exists, so \emph{revenue is understated}; and
because the marginal cycles the fee should have suppressed are still present,
\emph{throughput and cycle count are overstated}. Measured against a re-solved
run on a full synthetic year, a 5~EUR/MWh charging fee understates revenue by
1.0\% while overstating throughput by 13.3\%; at 25~EUR/MWh the figures are
50.7\% and 101.1\%. The asymmetry matters, because the augmentation trigger keys
off cycle count and therefore fires early long before the revenue figure looks
wrong. The error is exactly zero while the exemption holds.

\section{Reserve Products (Specified, Not Implemented)}

\textbf{Nothing in this section is implemented.} The parameters below are
present in \nolinkurl{scenarios/reference.yaml}, but \nolinkurl{config.py} defines no
FCR or aFRR dataclass and does not read them, and \nolinkurl{solve_window} builds
no reserve variables. What follows is the formulation these parameters imply,
recorded so the schema and the intended model can be checked against each other
before either is built.

\subsection{Additional sets and parameters}

Let $k \in \mathcal{K}$ index four-hour reserve blocks, with $\mathcal{T}_k$ the
model steps in block $k$ and $H^{\text{blk}}$ the block length in hours.

\small
\begin{tabular}{@{}l >{\raggedright\arraybackslash}p{6.0cm} >{\raggedright\arraybackslash}p{6.6cm}@{}}
\toprule
Symbol & Description & Source field \\
\midrule
$c^{\text{FCR}}$      & FCR capacity price (EUR/MW/h) & \nolinkurl{fcr.capacity_price_eur_per_mw_h} \\
$h^{\text{FCR}}$      & Energy headroom held per MW awarded (h) & \nolinkurl{fcr.reserve_energy_hours} \\
$R^{\min}$            & Minimum bid size (MW) & \nolinkurl{fcr.min_bid_mw}, \nolinkurl{afrr.min_bid_mw} \\
$c^{a\pm}$            & aFRR capacity price, positive direction shown; negative is the counterpart & \nolinkurl{afrr.positive.capacity_price_eur_per_mw_h} \\
$\pi^{a\pm}$          & aFRR energy price, likewise per direction & \nolinkurl{afrr.positive.energy_price_eur_per_mwh} \\
$\alpha^{\pm}$        & Expected activation rate, per direction & \nolinkurl{afrr.positive.activation_rate} \\
$h^{a}$               & aFRR energy headroom per MW awarded (h) & \nolinkurl{afrr.reserve_energy_hours} \\
\bottomrule
\end{tabular}
\normalsize

\subsection{Additional variables}

$R^{\text{FCR}}_{k} \ge 0$ is the symmetric FCR award in block $k$;
$R^{a+}_{k}, R^{a-}_{k} \ge 0$ are the directional aFRR awards. Binary
indicators $y^{\bullet}_{k} \in \{0,1\}$ carry the minimum bid size.

\subsection{Revenue terms}

The objective (2.1) gains capacity revenue and expected activation revenue:

\begin{equation}
+ \sum_{k \in \mathcal{K}} H^{\text{blk}}
  \Big[ c^{\text{FCR}} R^{\text{FCR}}_{k}
      + c^{a+} R^{a+}_{k} + c^{a-} R^{a-}_{k}
      + \alpha^{+} \pi^{a+} R^{a+}_{k}
      - \alpha^{-} \pi^{a-} R^{a-}_{k} \Big]
\end{equation}

where the sign on the negative-direction energy term reflects that
$\pi^{a-}$ is quoted as payment \emph{to} the provider for absorbing energy, so
a negative price is revenue. This convention must be pinned down against
settlement rules before implementation; it is stated here precisely because it
is the kind of sign that is easy to get backwards and hard to detect afterwards.

\subsection{R1 --- Minimum bid size}

\begin{equation}
R^{\min} y^{\bullet}_{k} \le R^{\bullet}_{k} \le \bar{P} \, y^{\bullet}_{k}
\qquad \forall k \in \mathcal{K}
\end{equation}

where $\bullet$ ranges over the three products. An award is either zero or at
least $R^{\min}$, which is a semi-continuous variable and the reason reserve
participation introduces binaries of its own.

\subsection{R2 --- Power headroom}

Capacity sold must be deliverable on top of energy dispatch:

\begin{align}
P^{d}_{t} + R^{\text{FCR}}_{k} + R^{a+}_{k} &\le \bar{P}
  & \forall k, \; \forall t \in \mathcal{T}_k \\
P^{c}_{t} + R^{\text{FCR}}_{k} + R^{a-}_{k} &\le \bar{P}
  & \forall k, \; \forall t \in \mathcal{T}_k
\end{align}

\subsection{R3 --- Energy headroom}

Capacity must also be sustainable for the prequalification duration, which is
what $h^{\text{FCR}}$ and $h^{a}$ encode:

\begin{align}
E_{t} - \big( h^{\text{FCR}} R^{\text{FCR}}_{k} + h^{a} R^{a+}_{k} \big)
  &\ge E^{\min} & \forall k, \; \forall t \in \mathcal{T}_k \\
E_{t} + \big( h^{\text{FCR}} R^{\text{FCR}}_{k} + h^{a} R^{a-}_{k} \big)
  &\le E^{\max} & \forall k, \; \forall t \in \mathcal{T}_k
\end{align}

where $h^{\text{FCR}} = 0.5$ reflects the German requirement that full FCR
activation be sustainable for 30 minutes in both directions. Understating it
permits the same MW to be sold to reserve and to arbitrage at once.

\subsection{Known gap: activation does not move the state of charge}

In the formulation above, aFRR energy enters revenue through the expected
activation rate $\alpha^{\pm}$ but does not appear in the state-of-charge
dynamics (2.4). Real activation moves energy and therefore constrains subsequent
dispatch. Closing that gap requires either activation as a deterministic energy
flow inside B2, or a stochastic formulation over activation scenarios. Until it
is closed, reserve participation will be overvalued, because the model collects
activation revenue without paying its state-of-charge consequence.

\section{Execution and Reported Quantities}

\subsection{Running a scenario}

One scenario file in, one revenue summary out:

\begin{verbatim}
python scripts/run_scenario.py scenarios/reference.yaml           # full year
python scripts/run_scenario.py scenarios/reference.yaml --days 7  # first 7 days
\end{verbatim}

\texttt{bess-stack-run} is the same code installed as a console script, and
appears once the package has been reinstalled, since the installer is what
generates it.

\subsection{Definition of the reported quantities}

The summary reports six quantities, defined here in terms of the decision
variables of Section 2.3 so that a number in the output can be traced to the
model that produced it.

\begin{align}
\text{gross revenue} &= \sum_{t \in \mathcal{T}} \pi_t \big( P^{d}_{t} - P^{c}_{t} \big) \Delta t \\
\text{energy discharged} &= \sum_{t \in \mathcal{T}} P^{d}_{t} \Delta t \\
\text{energy charged} &= \sum_{t \in \mathcal{T}} P^{c}_{t} \Delta t \\
\text{degradation cost} &= c^{\deg} \sum_{t \in \mathcal{T}} P^{d}_{t} \Delta t \\
\text{net revenue} &= \text{gross revenue} - \text{degradation cost} \\
\text{equivalent full cycles} &= \frac{1}{E^{\text{nom}}} \sum_{t \in \mathcal{T}} P^{d}_{t} \Delta t
\end{align}

Note that (4.1) is evaluated at the raw price $\pi_t$, whereas the objective
(2.1) is evaluated at the captured price $\pi_t \kappa$ net of fees. The
reported net revenue is therefore \emph{not} the objective value in general:
the two differ by

\begin{equation}
(1 - \kappa) \sum_{t} \pi_t \big( P^{d}_{t} - P^{c}_{t} \big) \Delta t
\; + \; \big( f + g^{d} \big) \sum_{t} P^{d}_{t} \Delta t
\; + \; \big( f + g^{c} \big) \sum_{t} P^{c}_{t} \Delta t .
\end{equation}

In the reference scenario $\kappa = 1$ and $f = g^{c} = g^{d} = 0$, so the
difference vanishes and the two coincide. That is a property of the reference
case rather than an identity, and it stops holding the moment a haircut or a fee
is set.

Energy discharged is the quantity degradation is charged on, and equivalent full
cycles is that same energy expressed in units of nominal capacity. The two carry
the same information and are reported together only because the cycle count is
the figure that feeds the augmentation trigger.

\subsection{Reference case}

Measured on the reference scenario over a full synthetic year, not illustrative:

\begin{tabular}{@{}lr@{}}
\toprule
Quantity & Value \\
\midrule
Gross revenue           & 473{,}374 EUR \\
Degradation cost        & 46{,}104 EUR \\
Net revenue             & 427{,}270 EUR \\
Energy discharged       & 11{,}526 MWh \\
Energy charged          & 13{,}565 MWh \\
Equivalent full cycles  & 576.3 \\
\bottomrule
\end{tabular}

The gap between energy charged and energy discharged is round-trip loss plus the
idle loss $L$ of (2.5), and is the physical content of the efficiency asymmetry
described under constraint group B2.

\subsection{What the reported figures exclude}

The summary carries four caveats. They are not decoration on the output; they
are the difference between what the scenario describes and what the model
computes, and each is a restatement in operational terms of a boundary already
drawn in this document.

\begin{itemize}
  \item \textbf{Day-ahead only.} \nolinkurl{revenue_streams.fcr} and
        \nolinkurl{revenue_streams.afrr} are enabled in the reference scenario,
        but Section 3 is not implemented, so every figure is the day-ahead leg
        alone and understates the stack the scenario describes.
  \item \textbf{No return metrics.} Nothing in \nolinkurl{finance/} computes the
        quantities \nolinkurl{outputs.metrics} names, so the summary stops at
        annual net revenue. That number is not a return: it is gross margin
        before capital cost, tax and financing.
  \item \textbf{Synthetic prices.} \nolinkurl{market.price_source} is
        \texttt{synthetic}, and the alternatives raise rather than silently
        degrade. The figures measure the price generator of
        \nolinkurl{data/prices.py}, not the German market.
  \item \textbf{Short runs are not annualised.} \texttt{-{}-days N} models the
        first $N$ days of the series, which begin in January. January carries
        the narrowest spreads in the synthetic year, so pro-rating a short run
        understates the year rather than approximating it. Seven days yield
        5{,}940 EUR net, which pro-rates to roughly 310{,}000 EUR against the
        427{,}270 EUR measured above --- an understatement of about 27\%.
\end{itemize}

Network charges are reported separately and are zero in every year of the
reference case, for the reason given in Section 2.7: the exemption of
\S118(6) EnWG runs to 2044 and the modelled life is 2025--2044, so no project
year is charged.

\end{document}
